% Imported from topics/video-generation/seedance-2.5/av-editing-spec.md
\chapter{音视频编辑状态与四类掩码}\label{seedance-m23}
\moduleinfo{M23}{候选设计与工程选型；组合效果待验证}{\hyperref[src-vace]{\texttt{vace}}、\hyperref[src-videocof]{\texttt{videocof}}、\hyperref[src-vincie-next]{\texttt{vincie-next}}、\hyperref[src-ltx-2]{\texttt{ltx-2}}、\hyperref[src-ltx-original-code]{\texttt{ltx-original-code}}、\hyperref[src-ltx-original-model]{\texttt{ltx-original-model}}}

接着把源条件扩到音视频：分别决定哪条流、哪个区间需要更新，最终在哪个表示域保持原内容。

同样叫局部编辑，改外套与改台词需要更新的模态并不相同。前者可能要求声音逐采样保持；后者则要改变声音并让嘴形跟上。因此先定义保留规则，再决定视频和声音中哪些位置进入去噪，不能从“局部”二字直接推出同一种输入。

这里仍在同一音视频主干上增加干净源条件和区域支路。严格外区由最终合成保持，几何不再对齐的任务另按语义检查。下面把这个候选过程逐步写出，尚未经过实际训练。

\section{一次编辑保存四类mask}\label{seedance-m23-ux4e00ux6b21ux7f16ux8f91ux4fddux5b58ux56dbux7c7bmask}

{\def\LTcaptype{table} % do not increment counter
\small\setlength{\tabcolsep}{3pt}\renewcommand{\arraystretch}{1.18}
\begin{longtable}[]{@{}
  >{\raggedright\arraybackslash}p{(\linewidth - 4\tabcolsep) * \real{0.1900}}
  >{\raggedright\arraybackslash}p{(\linewidth - 4\tabcolsep) * \real{0.3900}}
  >{\raggedright\arraybackslash}p{(\linewidth - 4\tabcolsep) * \real{0.4200}}@{}}
\toprule\noalign{}
\begin{minipage}[b]{\linewidth}\raggedright
名称
\end{minipage} & \begin{minipage}[b]{\linewidth}\raggedright
决定什么
\end{minipage} & \begin{minipage}[b]{\linewidth}\raggedright
典型形状
\end{minipage} \\
\midrule\noalign{}
\endhead
\bottomrule\noalign{}
\endlastfoot
条件可见性 & query允许读哪些源/参考/历史 & 逻辑关系或块稀疏结构 \\
去噪更新 & 采样器改变哪些目标位置 & 视频 \(B,N_v,1\)，音频 \(B,N_a,1\) \\
监督有效位 & 哪些预测计入训练损失 & 与目标token对齐 \\
最终合成alpha & 交付文件替换哪些像素/采样 & 像素帧网格或波形采样轴 \\
\end{longtable}
}

四种 mask 虽可从同一用户范围派生，作用却不相同。VAE 有时间感受野，模型工作区往往比最终替换区大；甚至可以先生成整个目标，再只合成一个小区域。共用一份不加说明的 mask，就把这些区别抹掉了。

\section{从时间戳到latent不是简单除八}\label{seedance-m23-ux4eceux65f6ux95f4ux6233ux5230latentux4e0dux662fux7b80ux5355ux9664ux516b}

用户区间使用共同媒体时钟的半开区间 \([a,b)\)。视频帧由PTS选择，音频按采样时间选择；VFR视频先建立真实重采样映射。

latent先保存原生位置网格边界，工作区选择与 \([a,b)\) 相交的token，再依据encoder、源塔与decoder的实际感受野扩张。位置网格只用于时间对应，不表示某个token只依赖或只影响这段像素。空间mask按主体跟踪传播，并考虑下采样与卷积覆盖；只取中心点可能漏掉边缘。

还要看解码端的依赖。原系列视频 encoder 时间因果，并不推出 decoder 不读未来；直接逐块解码，未必等于整段解码。因此首版先保存完整 latent、统一解码再合成。若以后分块，还需按实际感受野安排 halo 和边界处理。

\section{源条件保持干净，目标可以全部去噪}\label{seedance-m23-ux6e90ux6761ux4ef6ux4fddux6301ux5e72ux51c0ux76eeux6807ux53efux4ee5ux5168ux90e8ux53bbux566a}

源条件先拆待变与保留部分，各自编码；参考与目标也独立编码。首版在工作窗口内对目标完整去噪，再在像素端处理严格保持，避免先假设每个保留像素都对应一个完全独立的已知 latent。

已知区重噪覆盖是可选的第二种策略：

\[
z_{\sigma}^{\mathrm{known}}=(1-\sigma)z_0^{\mathrm{source}}+
\sigma\epsilon^{\mathrm{fixed}}.
\]

若采用已知区覆盖，每步要沿同一条固定噪声路径，训练中也见过相同方式。它和干净源条件不同，仍不能保证像素无损。边界像素共享 latent 时，可先让边界 token 进入工作区，最后按允许区域合成。

\section{先区分四种音轨任务}\label{seedance-m23-ux5148ux533aux5206ux56dbux79cdux97f3ux8f68ux4efbux52a1}

{\def\LTcaptype{table} % do not increment counter
\small\setlength{\tabcolsep}{3pt}\renewcommand{\arraystretch}{1.18}
\begin{longtable}[]{@{}
  >{\raggedright\arraybackslash}p{(\linewidth - 6\tabcolsep) * \real{0.1600}}
  >{\raggedright\arraybackslash}p{(\linewidth - 6\tabcolsep) * \real{0.2800}}
  >{\raggedright\arraybackslash}p{(\linewidth - 6\tabcolsep) * \real{0.2800}}
  >{\raggedright\arraybackslash}p{(\linewidth - 6\tabcolsep) * \real{0.2800}}@{}}
\toprule\noalign{}
\begin{minipage}[b]{\linewidth}\raggedright
编辑任务
\end{minipage} & \begin{minipage}[b]{\linewidth}\raggedright
视频目标
\end{minipage} & \begin{minipage}[b]{\linewidth}\raggedright
音频目标
\end{minipage} & \begin{minipage}[b]{\linewidth}\raggedright
交付策略
\end{minipage} \\
\midrule\noalign{}
\endhead
\bottomrule\noalign{}
\endlastfoot
换衣/换背景，声音保持 & 工作区生成 & 无；源音轨干净条件 & 外区原像素＋完整原音轨 \\
修改台词 & 嘴形及允许影响区域生成 & 指定台词区间生成 & 原音轨外区＋新台词区间 \\
改动作及对应音效 & 动作与接触相关区域生成 & 相应音效工作区生成 & 声画统一验收，允许必要上下文变化 \\
整体换风格或相机 & 语义约束下生成 & 按任务保持或重生 & 几何/身份/同步验收，不能要求全部外区逐像素不变 \\
\end{longtable}
}

台词任务还要给转录、speaker、时间预算和声音参考。新句子若长于原片段，就必须重写或明确允许延长；偷偷截掉字、任意拉伸声音，都没有满足原来的逐字与时间要求。

\section{固定音频与局部音频怎样参加去噪}\label{seedance-m23-ux56faux5b9aux97f3ux9891ux4e0eux5c40ux90e8ux97f3ux9891ux600eux6837ux53c2ux52a0ux53bbux566a}

固定音轨每步使用同一干净音频latent，噪声时间为0，声音流参与隐藏状态计算，音频损失与latent更新关闭。局部重生成时，仅活动区间加噪和更新，外部为干净条件。

候选分布为：

\[
\sigma_i^a=A_i\sigma,\quad
z_i^a=(1-\sigma_i^a)z_{0,i}^a+\sigma_i^a\epsilon_i^a.
\]

\(A_i\)表示音频工作区。它必须从codec覆盖区间派生，不能把每个40毫秒位置视为互不重叠的独立波形块。混响与卷积上下文也会跨过最终剪辑边界。

状态构造的机制代码见 \hyperref[code-condition-state]{condition\_state.py}，仅作说明。这里再强调一个依赖：固定输入音频 latent 后，AV 交互仍读取每步的视频，所以深层声音特征不固定，默认仍逐步重算。

\section{自动区域错误怎样回退}\label{seedance-m23-ux81eaux52a8ux533aux57dfux9519ux8befux600eux6837ux56deux9000}

无用户mask时，先用指令定位、检测与跟踪产生候选区域。区域ID、置信度和目标主体槽一起进入第二阶段。定位错误与生成错误分别标注，低置信目标可退回关键帧修正。

定位输出也有误差，生成阶段必须见过它。训练可混入可控腐蚀、膨胀、漂移以及真实预测区域，避免始终用完美 mask 训练。灰度区域生成和分割头是两种实现，首版先固定其一，使输出与监督含义对应。

\section{遮挡、阴影与运动边界怎样定义}\label{seedance-m23-ux906eux6321ux9634ux5f71ux4e0eux8fd0ux52a8ux8fb9ux754cux600eux6837ux5b9aux4e49}

例如手或道具挡住外套，mask 就还要记录衣服的可见部分、遮挡物归属和边界置信度。若遮挡物应保持，可单独复用原 alpha；若要生成新姿态，被遮住区域又可能必须重建。保持范围因此要随任务改变。

颜色变化若需要改变反光和阴影，应显式增加允许影响区域。若用户只批准外套像素，输出可以严格局部合成，但不能同时保证全局光照物理一致。

\section{交付合成的保证范围}\label{seedance-m23-ux4ea4ux4ed8ux5408ux6210ux7684ux4fddux8bc1ux8303ux56f4}

合成放在高质量工作表示上：视频按 alpha 取源帧或生成帧，音频按原采样率和时间取相应波形，之后统一编码混流。这里保持的是编码前外区；有损重编码会再次改变表示，所以不能连带保证文件字节或最终解码像素完全相同。

音频交界可以交叉淡化，淡化区必须包含在允许编辑时间内。声学连续性与逐采样外区保持冲突时，需要在任务规格中明确哪项优先，不能把淡化悄悄扩到禁止区间。

绿幕输出还需要独立前景alpha、边缘去色溢出和颜色管理。背景替换后的人物光照若需变化，应放入可编辑条件；单独输出一个绿色背景不是可靠matte。

\section{怎样判断改对而不是没改}\label{seedance-m23-ux600eux6837ux5224ux65adux6539ux5bf9ux800cux4e0dux662fux6ca1ux6539}

评价时把目标变化与外区保持放在一起，再看边界、身份、动作、台词、speaker 和声画偏差。只让保持误差变小，复制整段源视频也能做到，因此它不能单独证明编辑有效。

真实条件组织看 \href{https://github.com/ali-vilab/VACE}{VACE}，区域中间表示看 \href{https://github.com/knightyxp/VideoCoF}{VideoCoF}。已有LTX \href{https://github.com/Lightricks/LTX-2/blob/bd92a5f408eb8c99f8c05e9fe2fcaaec8b56b638/packages/ltx-trainer/src/ltx_trainer/training_strategies/video_to_video.py}{视频IC-LoRA策略} 是实现起点；上述音频状态、区域损失和最终合成需要另外接入。
